\section{Theoretical Background}
\label{sec:Theory}
In particle physics, an important concept is the cross section $\sigma$. 
The cross-section describes the probability for an interaction to happen. 
It can also be imagined as the effective target area for the projectile \cite{Povh1995}.
It depends on the projectile, its energy and the target. The unit of the cross section is $\SI{1}{\barn}=\SI{e-28}{\meter\squared}$.
For this experiment the more important question is: where does the scattered particle go?
This information is given by the differential cross section
\begin{equation}
    \frac{\text d \sigma}{\text d \Omega} = \frac{1}{\dot N_0 \, n x}\frac{\text d \dot N}{\text d \Omega} \, ,
    \label{eqn:diffsigma}
\end{equation}
which describes the fraction of particles scattered into the solid angle element $\text d \Omega$ at the angle $\theta$ \cite{Povh1995}.
Here, $\dot N_0$ is the rate of incident particles, $n$ the number density of the target nuclei and $x$ the thickness of the foil.
This holds as long as the foil is thin enough that every particle is scattered at most once.
For a nucleus with only a Coulomb potential, it is given by the Rutherford formula, which is discussed in the next section.
Since the Coulomb potential has an infinite range, the total cross section $\sigma$ diverges in this case, so only the differential cross section can be measured.

\subsection{Rutherford scattering}

The model tested by the Rutherford experiment is the so-called \enquote{plum pudding} or Thomson model, 
which describes the atom as a homogeneous distributed positive charge in which the electrons are embedded \cite{demtroeder3}.
In this model, only small deflection angles are possible.
However, the scattering experiment by Geiger and Marsden showed that this is not the case \cite{geiger1913}.
Most $\alpha$-particles passed the foil almost undeflected, but some were scattered by large angles, even above $90°$.
Rutherford's explanation is the nuclear model of the atom, so that
the positive charge and almost the entire mass are concentrated in a very small nucleus, which is surrounded by the electrons \cite{rutherford1911}.
\\
For the scattering, Rutherford derived a formula that describes the scattering. 
The assumptions of that are that the projectile and target are only point charges, the whole 
force on the projectile is only the Coulomb force; the nuclei of the target are at rest and cannot move. 
It is classical and not relativistic, and there is only one scattering, and no energy is lost in the scattering. 
This results in
\begin{equation}
    \label{eqn:rut}
    \frac{\text d \sigma}{\text d \Omega} = \left(\frac{1}{4\pi\epsilon_0} \right)^2\left(\frac{zZe^2}{4E} \right)^2\frac{1}{\sin^4(\theta/2)}\ .
\end{equation}
There is $z$, the nuclear charge number of the projectile, and $Z$, the one of the target. 
$E$ is the kinetic energy of the projectile and $\theta$ the angle in which the projectile is being scattered \cite{Povh1995}.

\subsection{Alpha decay and the \texorpdfstring{\ce{^{241}Am}}{241Am} source}

The alpha decay can be written as $\ce{^A_ZX} \rightarrow \ce{^{A-4}_{Z-2}Y}+ \ce{^4_2He^{2+}}$ where the $\ce{^4_2He^{2+}}$ is the $\alpha$ particle itself.
Americium-241 is used as the $\alpha$ emitter.
The half-life is $T_{1/2}= \SI{432.2}{a}$ with an activity of $A = \SI{330}{kBq}$ measured in October 1994, so the current activity is somewhat lower than that \cite{Durchfuhrung}.
The term diagram of \ce{^{241}Am} is given in \autoref{fig:td}.
\begin{figure}[htbp]
    \centering
    \includegraphics[width=0.75\textwidth]{v16/content/pics/The-Decay-Scheme-of-Am-241.pdf}
    \caption{Decay scheme of \ce{^{241}Am}. Picture from \cite{inproceedings}}
    \label{fig:td}
\end{figure}
The most important decay is the one with an energy of $\SI{5.486}{MeV}$. Also given are some $\gamma$ decays with energies in the keV range. 
So it is possible to differentiate between the $\alpha$ and the $\gamma$ decays \cite{inproceedings}. 

\subsection{Bethe-Bloch equation}

Charged particles lose energy when passing through matter. 
This is caused by many small Coulomb collisions with atomic electrons, which ionise and excite the atoms \cite{leo}.
This energy loss per distance is described by the Bethe-Bloch equation 
\begin{equation}
    -\frac{\text{d}E}{\text{d}x} = \frac{4 \pi e^4 z^2 N Z}{m_0 v^2 \left(4 \pi \epsilon_0\right)^2}
    \left[ \ln\left(\frac{2 m_0 v^2}{I}\right) - \ln\left(1 - \beta^2\right) - \beta^2 \right] \ .
    \label{eq:bethe}
\end{equation}
Here $z$ and $v$ are the charge number and velocity of the incident particle, $\beta = v/c$,
$N$ is the number density of the target atoms, $Z$ their atomic number, $m_0$ the electron mass
and $I$ the mean ionization energy of the target material.
An important dependency is the $1/\beta^2$ which says that slower projectiles are depositing more energy. 
Through that, the deposit energy peaks just before it stops, which is called the "Bragg-Peak". 
The range of a charged particle is calculated through
\begin{equation}
    R = \int_0^{E_0}\left(- \frac{\text d E}{\text d x }\right)^{-1}\text d E \ .
\end{equation}

The Bethe-Bloch equation assumes that the projectile is heavy compared to an electron, so it is hardly deflected.
In addition, the electrons are assumed to be quasi-free and at rest, and the projectile has to be much faster than the bound electrons.
For this experiment it is important that the energy loss per distance is dependent on the pressure of the air, so 
\begin{equation}
    -\frac{\text{d}E}{\text d x} \propto p \ .
\end{equation}
For the measurement of the thickness of the foil, this relation is used. 
It is used by using the linear dependence of the energy loss on the pressure. 
Because the rest energy can be measured through the surface barrier detector, which will be described in the following.


\subsection{Surface barrier detector}

The detector is based on silicon, which has a band gap of $E_g \approx \SI{1.12}{eV}$ at room temperature.
Because of that only a few charge carriers are thermally excited.
If an n-doped and a p-doped region meet, the majority carriers diffuse across the junction and recombine.
This creates a depletion zone without free charge carriers, in which the fixed ionised dopant atoms build up an electric field.
This depletion zone is the sensitive volume of the detector, because the charge created there by the $\alpha$ particle is separated by the field and collected.
\\
The surface barrier detector consists of an n-type silicon wafer with a very thin layer of gold on the front side.
The gold forms a Schottky contact with the silicon, which acts like a pn junction, so the depletion zone is directly below the surface.
Because the gold layer is so thin, the $\alpha$ particles only lose a few $\si{keV}$ before they reach the sensitive volume.
\\
The detector is operated in reverse bias, which widens the depletion zone.
Its depth is given by
\begin{equation}
    w = \sqrt{\frac{2\varepsilon\left(U + U_\text{bi}\right)}{e N_D}} \ ,
\end{equation}
with the bias voltage $U$, the built-in voltage $U_\text{bi}$ and the donor concentration $N_D$.
The depth $w$ has to be larger than the range of the $\alpha$ particles in silicon, which is about $\SI{25}{\micro\meter}$ to $\SI{30}{\micro\meter}$ at $\SI{5.5}{MeV}$.
Only then the particles deposit their whole energy in the sensitive volume and the measured signal is proportional to their energy.
The $\alpha$ particles creates electron-hole pairs, one per $\approx \SI{3.6}{eV}$. The collected charge is therefore proportional to the deposited energy.
A disadvantage of the surface barrier detector is its sensitivity to light, which is why the chamber has to be covered so no light is hitting the detector \cite{leo}.
\subsection{Error Calculation}
% Error Calculation
The error calculation in this protocoll is being done using the gaussian error propagation
\begin{equation}
    \sigma_f = \sqrt{\left( \frac{\partial f}{\partial x_1} \sigma_{x_1} \right)^2 + \left( \frac{\partial f}{\partial x_2} \sigma_{x_2} \right)^2 + ...} \, ,
\end{equation}
as well as the formulas 
\begin{equation}
    \overline{x} = \frac{1}{n} \sum^n_{i=1}x_i \, ,
    \label{eq:mittelwert}
\end{equation}
for the mean and 
\begin{equation}
    \sigma_{\overline{x}} = \frac{\sigma_x}{\sqrt{n}} = \sqrt{\frac{1}{n(n-1)}  \sum^n_{i=1} \left( x_i - \overline{x}\right)^2 } \, ,
    \label{eq:abweichung}
\end{equation}
for the standard deviation. 
To simplify these calculations, the Python-package \textit{uncertainties}~\cite{uncertainties} is being used. 

\clearpage